\section{Signed logarithmic energy and smoothing}
\label{sec:v-analytic}

For each fixed $q\in\{2,3,4,5,6,7\}$, directed interval checks on the printed
reference and the smoothing estimate below give an all-$m$ determinant
bound.  The interval checks certify the energy rate; the smoothing lemma
supplies the estimate for every $m=4n\ge64$.

Throughout, $q$ is a fixed integer with $2\le q\le7$, $a=-q$,
$m=4n$.  The three local shape parameters are set to $s=k=h=n$, so that
$L=m+h=5m/4$.  The three masses are
$(M_x,M_v,M_z)=(1,5/4,1)$.  Write
\[
 I(\sigma,\eta)=\iint\log|u-w|\,d\sigma(u)d\eta(w),\qquad
 D_q(\sigma,\eta)=\iint\log|q+uv|\,d\sigma(u)d\eta(v),
\]
and
\begin{equation}\label{eq:v-energy}
 \mathcal Q_q(\mu)=I(\mu_x,\mu_x)+I(\mu_v,\mu_v)+I(\mu_z,\mu_z)
       +I(\mu_v,\mu_z)-D_q(\mu_x,\mu_v).
\end{equation}
The supports used for the continuous problem are
\[
 \mathcal K_x=\mathcal K_v=[0,1],\qquad
 \mathcal K_z(q)=[-q,-r]\cup[r,1]\cup\{|z|=r\},\qquad r=3/5.
\]
For a reference triple $\nu=(\nu_x,\nu_v,\nu_z)$ of these masses, put
\begin{align*}
 U_i(w)&=\int\log|w-t|\,d\nu_i(t),\\
 D_v(w)&=\int\log|q+wt|\,d\nu_v(t),\\
 D_x(w)&=\int\log|q+tw|\,d\nu_x(t).
\end{align*}
and define the three fields
\begin{align}
 \widehat F_x(w)&=2U_x(w)-D_v(w),\notag\\
 \widehat F_v(w)&=2U_v(w)+U_z(w)-D_x(w),\label{eq:v-q-fields}\\
 \widehat F_z(w)&=2U_z(w)+U_v(w),\notag
\end{align}
For the full logarithmic energy we add the one-body terms separately:
\begin{align}
 F_x(w)&=\widehat F_x(w)+\tfrac34\log|w|,\notag\\
 F_v(w)&=\widehat F_v(w)+\tfrac14\log|w|,\label{eq:v-fields}\\
 F_z(w)&=\widehat F_z(w)-\tfrac54\log|w|.\notag
\end{align}
The constant field from the determinant prefactor is
\begin{equation}\label{eq:v-C0}
 C_0(q)=\tfrac32\log q.
\end{equation}

We first identify the determinant at the dilogarithm value. Put
\[
 M_{ij}(q)=\int_0^1\!\int_0^1\frac{x^i y^j}{q+xy}\,dx\,dy,
 \qquad \beta_q=\operatorname{Li}_2(-1/q),
 \qquad \lambda_q=\log(1+1/q).
\]
For $j\ge0$ set the two moment sequences
\[
 m_j=\int_0^1\frac{u^j}{q+u}\,du,
 \qquad n_j=\int_0^1\frac{u^j\log(1/u)}{q+u}\,du.
\]
Polynomial division gives, for $j\ge1$,
\[
 m_j=-q m_{j-1}+\frac1j,\qquad
 n_j=-q n_{j-1}+\frac1{j^2},
\]
while $m_0=\lambda_q$ and $n_0=-\beta_q$. Hence, with the
polynomials $t_j,b_j$ of Section~\ref{sec:v-arithmetic},
\[
 m_j=(-q)^j\lambda_q+t_j,\qquad n_j=b_j-(-q)^j\beta_q.
\]
The substitution $u=xy$ yields
\[
 M_{ij}(q)=\frac{m_i-m_j}{j-i}\quad(i\ne j),
 \qquad M_{ii}(q)=n_i,
\]
and therefore
\begin{equation}\label{eq:v-moment-identity}
 M(q)=B(-q)+\lambda_q C(-q)-\beta_q A(-q).
\end{equation}
Adding $\lambda_q$ times each $C$ row to the corresponding row of
$B(-q)-\beta_qA(-q)$ preserves the determinant. Thus the determinant of
Section~\ref{sec:v-arithmetic}, evaluated at $T=\beta_q$,
is exactly the determinant with middle block $M(q)$.

\subsection{The path and the signed energy}

Let $\gamma_+(t)=re^{it}$ for $0\le t\le\pi$ and
$\gamma_-(t)=re^{it}$ for $\pi\le t\le2\pi$.  Orient the two real
segments increasingly and set
\begin{equation}\label{eq:v-path}
 \Gamma_q=[-q,-r]+[r,1]+(c_q-1)\gamma_++c_q\gamma_-,
 \qquad c_q=\frac{\log q+i\pi}{2\pi i}.
\end{equation}
For every integer $d$,
\begin{equation}\label{eq:v-laurent}
 \int_{\Gamma_q}z^{d-1}\,dz=
 \begin{cases}(1-(-q)^d)/d,&d\ne0,\\0,&d=0.\end{cases}
\end{equation}
For $d\ne0$, integrate $z^{d-1}$ on each piece using the primitive
$z^d/d$; the two circular half-arcs contribute exactly the terms needed
to cancel the intermediate endpoints $\pm r$.  For $d=0$, the real pieces
give $-\log q$, while the circular pieces give
$i\pi(2c_q-1)=\log q$.  This also fixes the branch convention at the
zero Laurent moment.

The same Andreief identity~\cite{Andreief,Forrester} and Cauchy--Vandermonde
expansions~\cite{Krattenthaler} as for the determinant
in Section~\ref{sec:v-arithmetic}, now using \eqref{eq:v-laurent}, give
\begin{equation}\label{eq:v-integral}
 G_m(\operatorname{Li}_2(-1/q);-q)
 =\frac{q^{3m^2/2-m}}{(m!)^2L!}
   \int_{[0,1]^m}\int_{[0,1]^L}\int_{\Gamma_q^m}
       \Phi_q(x,v,z)\,dz\,dv\,dx,
\end{equation}
where
\[
 \Phi_q=\frac{\Delta(x)^2\Delta(v)^2\Delta(z)^2
       \prod_{d=1}^m\prod_{b=1}^L(z_d-v_b)
       \prod_{i=1}^m x_i^{3m/4}
       \prod_{b=1}^L v_b^{m/4}
       \prod_{d=1}^m z_d^{-5m/4}}
       {\prod_{i=1}^m\prod_{b=1}^L(q+x_iv_b)}.
\]

\begin{proof}[Proof of the integral representation]
By \eqref{eq:v-laurent},
\(C_{ij}=(-q)^i\int_{\Gamma_q}z^{j-i-1}\,dz\).  Put
\(f_i(v)=\int_0^1x^i/(q+xv)\,dx\).  Then the lower rational block and the
V block are respectively
\[
 M_{ij}=\int_0^1f_i(v)v^j\,dv,\qquad
 V_{\ell j}=\int_0^1v^{\ell-1}v^j\,dv.
\]
For integrable families, Andreief's identity~\cite{Andreief,Forrester} is
\[
 \det\left(\int f_i(t)g_j(t)\,dt\right)_{i,j=1}^N
 =\frac1{N!}\int
   \det(f_i(t_b))_{i,b=1}^N\det(g_j(t_b))_{b,j=1}^N
   \prod_{b=1}^Ndt_b.
\]
Apply it successively to the $L$ lower/V rows, the $m$ contour rows, and
the $m$ functions $f_i$.  These three applications account respectively
for $1/L!$, $1/m!$, and $1/m!$:
\[
 G_m=\frac1{(m!)^2L!}\int D_zD_{zv}D_xD_{xv}\,dx\,dv\,dz.
\]
Here $e_\Delta=m(m-1)/2$ and, with
\(\mathbf w=(z_1,\ldots,z_m,v_1,\ldots,v_L)\), the four determinants are
\begin{align*}
 D_z&=\det\bigl((-q)^iz_d^{-i-1}\bigr)_{i\in I,d}
   =(-1)^{e_\Delta}(-q)^{\sum_{i\in I}i}
     \Delta(z)\prod_dz_d^{-3m/2},\\
 D_{zv}&=\det(\mathbf w_b^j)_{b,j\in J}
   =(-1)^{mL}\Delta(z)\Delta(v)
     \prod_dz_d^{m/4}\prod_bv_b^{m/4}\prod_{d,b}(z_d-v_b),\\
 D_x&=\det(x_{i'}^i)_{i\in I,i'}
   =\Delta(x)\prod_{i'}x_{i'}^{m/2},\\
 D_{xv}&=\det\begin{pmatrix}
 ((q+x_{i'}v_b)^{-1})_{i',b}\\(v_b^{\ell-1})_{\ell,b}
 \end{pmatrix}
   =\frac{(-1)^{e_\Delta}q^{e_\Delta}\Delta(x)\Delta(v)\prod_{i'}x_{i'}^{m/4}}
          {\prod_{i',b}(q+x_{i'}v_b)}.
\end{align*}
The last equality is the Cauchy--Vandermonde identity~\cite{Krattenthaler}.
For completeness,
for $t_{i'}=q/x_{i'}$ it follows from
\[
 \det\begin{pmatrix}((t_{i'}+v_b)^{-1})_{i',b}\\
 (v_b^{\ell-1})_{\ell,b}\end{pmatrix}
 =\frac{\Delta(t)\Delta(v)}{\prod_{i',b}(t_{i'}+v_b)},
\]
after multiplying the first $m$ rows by $x_{i'}^{-1}$ and using
\(\Delta(t)=(-1)^{e_\Delta}q^{e_\Delta}\Delta(x)\prod_{i'}x_{i'}^{1-m}\).
Here is a proof of the identity, for every $m,h\ge0$ with $L=m+h$.
Treat the $t_i$ and $v_b$ as independent indeterminates.  The case $m=0$
is the Vandermonde determinant.  Assume the formula for $m-1$ and all
$h\ge0$, and regard both sides as rational functions of $v_L$.  Expanding
the determinant along its last column, its residue at $v_L=-t_i$ is
$(-1)^{i+L}$ times the complementary minor.  By induction this residue is
\[
 (-1)^{i+L}
 \frac{\Delta(\widehat t_i)\Delta(v_1,\ldots,v_{L-1})}
      {\prod_{j\ne i}\prod_{b<L}(t_j+v_b)}.
\]
The proposed quotient has the same residue: substitute $v_L=-t_i$,
cancel the factors $t_j-t_i$ and $t_i+v_b$, and use
$\Delta(t)=(-1)^{i-1}\Delta(\widehat t_i)
\prod_{j\ne i}(t_j-t_i)$.  Hence the difference has no finite poles.
After cancellation of the denominator $\prod_i(t_i+v_L)$, the difference
is a polynomial in $v_L$ of degree at most $h-1$.  It vanishes at each of
the $L-1$ values $v_L=v_b$ ($b<L$), since both sides have the factor
$\Delta(v)$.  For $m\ge1$, $L-1=m+h-1\ge h$, so these roots force the
polynomial to be zero.  When $h=0$, the difference tends to zero at
infinity instead.  This proves the Cauchy--Vandermonde identity.

Multiplying the four displayed factors gives precisely $\Phi_q$: the
$m/2$ from $D_x$ and the $m/4$ from $D_{xv}$ combine to the displayed
$3m/4$ power of each $x_{i'}$.  The sign is
$(-1)^{\sum_{i\in I}i+2e_\Delta+mL}=1$: here
$\sum_{i\in I}i=2n(8n-1)$ and $mL=20n^2$ are even.  The q-power is
\[
 \sum_{i\in I}i+e_\Delta
 =m(2n)+m(m-1)
 =\frac{3m^2}{2}-m,
 \qquad (m=4n,\ L=5n).
\]
The factors $q+x_iv_b$ are at least
$q$, while $|z|\ge r$ on $\Gamma_q$; all three integrations therefore
converge absolutely.  The $m!$, $L!$, and $m!$ factors are exactly the three
Andreief symmetrizations, proving \eqref{eq:v-integral}.
\end{proof}
For triples of positive measures of the stated masses, define
\begin{equation}\label{eq:v-full-energy}
 \mathcal E_q(\mu)=C_0(q)+\mathcal Q_q(\mu)
 +\tfrac34\int\log|x|\,d\mu_x
 +\tfrac14\int\log|v|\,d\mu_v
 -\tfrac54\int\log|z|\,d\mu_z.
\end{equation}

\begin{lemma}[signed-energy concavity for fixed $q$]
\label{lem:v-concavity}
Let $\mu,\nu$ be triples of positive measures of masses
$(1,5/4,1)$, supported in compact subsets of the indicated supports or
their sufficiently small tubular neighborhoods.  Assume that the pairings
in $\mathcal Q_q$ are finite real numbers.  No absolute logarithmic moment
assumption is imposed on a competitor's $x$- or $v$-component.  Then
\[
 \mathcal Q_q(\mu)\le-\mathcal Q_q(\nu)
       +\int \widehat F_x\,d\mu_x+\int \widehat F_v\,d\mu_v
       +\int \widehat F_z\,d\mu_z,
\]
where the hatted fields are formed from $\nu$.  Consequently the same
inequality for the full energy is obtained by adding the displayed one-body
terms to the hatted fields, i.e. by using $F_x,F_v,F_z$ in the field bounds.
\end{lemma}

\begin{proof}
By the shared quadratic-duality identity in
Section~\ref{sec:certificate-algorithm}, it suffices to prove that the
zero-mass quadratic remainder $\mathcal Q_q(\mu-\nu)$ is nonpositive.
First suppose that all $v$-supports avoid zero.  For a signed measure
$\sigma$ of mass zero, write $\|\sigma\|^2=-I(\sigma,\sigma)\ge0$;
this is the zero-mass logarithmic-energy positivity lemma.  Put
$R_q(v)=-q/v$.  The exact identities
\begin{align*}
 \log|q+xv|&=\log|x-R_q(v)|+\log|v|,\\
 \log|R_q(v)-R_q(w)|&=\log q+\log|v-w|-\log|v|-\log|w|
\end{align*}
Integrating against zero-mass measures cancels all one-variable terms and
gives
\begin{align*}
 D_q(\sigma_x,\sigma_v)&=I(\sigma_x,(R_q)_*\sigma_v),\\
 I((R_q)_*\sigma_v,(R_q)_*\sigma_v)&=I(\sigma_v,\sigma_v).
\end{align*}
Let $\sigma_i=\mu_i-\nu_i$, and set
$\sigma'_x=\sigma_x$, $\sigma'_v=(R_q)_*\sigma_v$,
$\sigma'_z=-(R_q)_*\sigma_z$.  Then
\[
 \mathcal Q_q(\sigma)
 =-\|\sigma'_x\|^2-\|\sigma'_v\|^2-\|\sigma'_z\|^2
   +\langle\sigma'_x,\sigma'_v\rangle
   +\langle\sigma'_v,\sigma'_z\rangle,
 \qquad \langle\alpha,\beta\rangle=-I(\alpha,\beta).
\]
The right side equals
\[
 -\tfrac12\bigl(\|\sigma'_x-\sigma'_v\|^2
              +\|\sigma'_v-\sigma'_z\|^2
              +\|\sigma'_x\|^2+\|\sigma'_z\|^2\bigr)\le0.
\]
The shared identity gives the asserted quadratic inequality.  For a
competitor $\mu_v$ whose support meets zero, fix a smooth probability
measure $\rho$ supported in $[3/4,7/8]$.  For $0<\delta<1/4$, let
$\lambda_\delta=\mu_v|_{\{|v|<\delta\}}$, let
$e_\delta=\lambda_\delta(\mathbb C)$, and set
\[
 \mu_{v,\delta}=\mu_v-\lambda_\delta+e_\delta\rho.
\]
Finite logarithmic self-energy excludes an atom at zero, so
$e_\delta\to0$ and $\mu_{v,\delta}\Rightarrow\mu_v$.  Absolute
integrability of the logarithmic pairings in $\mathcal Q_q$ implies that
the integrals over the deleted sets tend to zero.  The replacement terms
involving $e_\delta\rho$ tend to zero as well: $\rho$ is separated from
$\{|v|<\delta\}$, and its logarithmic potential is bounded on the compact
supports.  The kernel $\log|q+uv|$ is bounded and continuous on the stated
tubes because $q\ge2$ and $|u|,|v|\le1+o(1)$.  Hence
$\mathcal Q_q(\mu_{x},\mu_{v,\delta},\mu_z)\to\mathcal Q_q(\mu)$.
The reference fields $\widehat F_i$ are continuous and bounded on these
compact sets, so their integrals also converge.  Applying the already
proved inequality to $\mu_{v,\delta}$ and taking the limit proves the
quadratic inequality without requiring $\int|\log|v||\,d\mu_v<\infty$.

The one-body integrals in $\mathcal E_q$ are then interpreted as extended
real integrals.  On the compact $x$- and $v$-supports, $\log|w|$ is bounded
above and can have integral $-\infty$; on the $z$-support, $|z|\ge r-o(1)$,
so the $z$-term is finite.  Adding these linear terms to the quadratic
inequality is therefore valid in the extended reals.  In particular, when
the full energy is $-\infty$ the asserted upper bound is automatic.  This
argument uses no implication from finite self-energy to a finite absolute
logarithmic moment.
\end{proof}

Each printed reference has finitely many positive-length rational interval
rows, compact $z$ support in $[-q,-3/5]$, and Poisson centres with
$|c|<3/5$; the $x$ and $v$ chains start at the positive endpoints listed in
the data section and end at $1$.  Hence all six reference supports are
compact and avoid zero.  Interval densities are bounded by the finite
maximum of their rational mass/length ratios, while a Poisson row of mass
$w$ has density
\[
 \frac{w}{2\pi}\frac{r^2-|c|^2}{|re^{i\theta}-c|^2}
 \le\frac{w(r+|c|)}{2\pi r(r-|c|)},\qquad r=3/5.
\]
Thus the printed references have finite logarithmic moments directly from
their rational supports, without any general self-energy implication.

Combining the lemma with the one-body terms gives the explicit full-energy
upper bound used below:
\begin{equation}\label{eq:v-full-energy-dual}
 \mathcal E_q(\mu)
 \le C_0(q)-\mathcal Q_q(\nu)
      +\int F_x\,d\mu_x+\int F_v\,d\mu_v+\int F_z\,d\mu_z.
\end{equation}
Thus, if the full fields are bounded above by $K_{x,q},K_{v,q},K_{z,q}$
on their respective admissible supports, then
\begin{equation}\label{eq:v-energy-bound}
 \mathcal E_q(\mu)\le -\bigl(\mathcal Q_q(\nu)-C_0(q)-K_{x,q}
                         -\tfrac54K_{v,q}-K_{z,q}\bigr).
\end{equation}

The inversion is valid uniformly for the fixed bases considered here:
on $|v|\le1+3/4096$,
$|R_q'(v)|=q/|v|^2>1$ for $q\ge2$.  Thus a bounded arclength density
pushes forward to a bounded arclength density on each inverted segment or
circle.  No $q=2$ numerical identity is used in the concavity argument.

\subsection{Directed fixed-base certificates}

For a reference triple $\nu_q$, suppose its energy lower endpoint is
$Q_{q,-}$ and its support-wide field upper bounds are
$F_x\le K_{x,q}$, $F_v\le K_{v,q}$, $F_z\le K_{z,q}$.  Lemma
\ref{lem:v-concavity} and \eqref{eq:v-full-energy} give the quadratic
rate
\begin{equation}\label{eq:v-alpha}
 \alpha_q=Q_{q,-}-\tfrac32\log q-K_{x,q}
                         -\tfrac54K_{v,q}-K_{z,q}.
\end{equation}
The following are the directed interval-bound results.  The field thresholds
are exact terminating decimals.  The displayed $\alpha$ values are strict
lower bounds rounded down, and each $\tau$ is an outward-rounded upper bound.

\begin{table}[ht]
\centering\normalsize
\renewcommand{\arraystretch}{1.12}
\caption{Fixed-base directed certificates for $a=-q$.}
\label{tab:v-certificates}
\setlength{\tabcolsep}{2pt}
\begin{tabular}{@{}c c c c c c c c@{}}
\hline
$q$&rows&$K_x$&$K_v$&$K_z$&$\alpha_q$&$\tau_q$&margin\\
\hline
2&$41/48/62/149$&$-4.621$&$-4.475$&$-0.085$&$>4.34875$&$<4.31765$\footnotemark&$>0.03110$\\
3&$40/48/62/143$&$-5.076$&$-4.632$&$0.388$&$>3.96707$&$<3.86995$&$>0.09712$\\
4&$40/48/62/141$&$-5.408$&$-4.727$&$0.779$&$>3.68211$&$<3.55230$&$>0.12981$\\
5&$39/48/62/137$&$-5.670$&$-4.791$&$1.108002$&$>3.45240$&$<3.30591$&$>0.14649$\\
6&$41/47/61/135$&$-5.886$&$-4.839$&$1.391$&$>3.26092$&$<3.10460$&$>0.15632$\\
7&$39/47/62/131$&$-6.071$&$-4.876$&$1.639$&$>3.09651$&$<2.93439$&$>0.16212$\\
\hline
\end{tabular}
\end{table}
\footnotetext{The exact rate is
\(\tau_2=A_{40}-(53/48)\log2=4.317643952609845\ldots\); the displayed
decimal is a directed upper bound.}

The energy lower endpoints used in \eqref{eq:v-alpha} are printed separately
so that the energy and field slacks can be checked independently:
\begin{table}[ht]
\centering
\renewcommand{\arraystretch}{1.12}
\caption{Directed lower endpoints for the reference energies.}
\label{tab:v-energy-lower}
\begin{tabular}{@{}c c@{}}
\hline
$q$ & $Q_{q,-}$\\
\hline
2&$>-4.911272$\\
3&$>-4.863007$\\
4&$>-4.776198$\\
5&$>-4.684191$\\
6&$>-4.595191$\\
7&$>-4.511619$\\
\hline
\end{tabular}
\end{table}

For each $q=2,\ldots,7$, put $\delta=1/250$.  The boxes
$0\le x,v\le\delta$ are discharged by the endpoint bounds
\eqref{eq:v-endpoint}.  On $[\delta,1]$ and on the two real $z$ intervals,
including the complete negative interval $[-q,-3/5]$, the subdivision
recursively bisects each rational midpoint.  A cell is accepted exactly when
its directed upper field interval is finite and strictly below the displayed
threshold; otherwise both children are visited, in left-first
order, with maximum depth $30$.  Circle cells use the same rule in the
parameter $t\in[0,1]$, and the endpoints $t=0,1$ are checked directly.
This is the complete partition rule.  The six references contain $1733$ rows
in all, and their rows are printed in Section~\ref{sec:v-reference-data}.
The smoothing lemma below supplies the complex-tube transfer and the all-$m$
estimate.

\subsection{Reference densities and directed endpoint disks}
\label{subsec:v-allm-data}

All constants below are computed from the signed decimal tables in
Section~\ref{sec:v-reference-data}.  For
each interval row $[a,b]$ of mass $w$, put $H=w/(b-a)$.  Let $H_{x,q}$,
$H_{v,q}$ and $H_{z,q}$ be the largest such value in the $x$, $v$ and
negative-real $z$ chains.  For a Poisson row $(c,w)$ write
\[
 \rho_c=2^{-120}\left\lceil2^{120}|c|\right\rceil,
 \qquad
 P_q=\sum_{(c,w)}\frac{w(r+\rho_c)}{r(r-\rho_c)}.
\]
The ceiling is evaluated by integer square root on the printed rational
coordinates, and the construction checks $\rho_c<r$ row by row.  Since the
Poisson density on $|z|=r$ is at most
$w(r+\rho_c)/(2\pi r(r-\rho_c))$ per unit arclength, and a circle meets a
disk of radius $t$ in arclength at most $2\pi t$, the reference measures
satisfy $\nu_i(B(w,t))\le A_{i,q}t$, where
\[
 A_{x,q}=2H_{x,q},\qquad A_{v,q}=2H_{v,q},\qquad
 A_{z,q}=2H_{z,q}+P_q.
\]
Define
\[
 B_{i,q}=6A_{i,q}+2M_i,\qquad
 \Sigma_q=2B_{x,q}+\frac72B_{v,q}+\frac{13}{4}B_{z,q},
 \qquad \mathcal C_q=10+3\Sigma_q.
\]
In the table, $\Sigma_q$ is rounded upward to $10^{-3}$ and this rounded
rational value is used in $\mathcal C_q$.  The path variation is
\[
 \mathcal T_q=q+1-2r+r\sqrt{\pi^2+\log^2q}.
\]
The $E_{x,q}$ and $E_{v,q}$ columns are rounded upward to $10^{-4}$;
the $\mathcal T_q$ column is rounded upward to $10^{-6}$.

For $\delta=1/250$, put $P(t)=t(\log t-1)$ and define the exact real
integrals
\[
 J_{x,q}=\int\log(t+\delta)\,d\nu_x(t),\quad
 J_{v,q}=\int\log(t+\delta)\,d\nu_v(t),\quad
 J_{z,q}=\int\log(|z|+\delta)\,d\nu_z(z).
\]
An interval contributes $w(P(b+\delta)-P(a+\delta))/(b-a)$; for a
negative $z$ interval use endpoints $-a,-b$ in increasing order.  The
Poisson rows together contribute their total mass times $\log(r+\delta)$.
For $|w_0|\le\delta$, the triangle inequality and the masses give
\begin{align}
 F_x(w_0)&\le E_{x,q}:=2J_{x,q}-\frac54\log(q-\delta)
                      +\frac34\log\delta,\notag\\
 F_v(w_0)&\le E_{v,q}:=2J_{v,q}+J_{z,q}-\log(q-\delta)
                      +\frac14\log\delta.\label{eq:v-endpoint}
\end{align}
The endpoint bounds are directed upper endpoints rounded upward at
$10^{-4}$.  The same table gives $\Sigma_q$ rounded upward at $10^{-3}$ and
$\mathcal T_q$ rounded upward at $10^{-6}$; throughout the proof set
$\mathcal C_q=10+3\Sigma_q$ using the printed $\Sigma_q$.
\begin{table}[ht]
\centering\normalsize
\renewcommand{\arraystretch}{1.12}
\caption{Endpoint and all-$m$ transfer data recomputed from the printed tables.}
\label{tab:v-transfer-constants}
\begin{tabular}{@{}crrrr@{}}
\hline
$q$&$E_{x,q}$&$E_{v,q}$&$\Sigma_q$ (upper)&$\mathcal T_q$ (upper)\\
\hline
2&$-6.2821$&$-4.4935$&$14092.327$&$3.730292$\\
3&$-6.7656$&$-4.6587$&$13219.499$&$4.796889$\\
4&$-7.1126$&$-4.7555$&$12702.122$&$5.860319$\\
5&$-7.3835$&$-4.8208$&$12351.983$&$6.917916$\\
6&$-7.6061$&$-4.8688$&$12095.881$&$7.969979$\\
7&$-7.7947$&$-4.9057$&$11898.583$&$9.017257$\\
\hline
\end{tabular}
\end{table}

The endpoint entries satisfy $E_{x,q}<K_{x,q}$ and $E_{v,q}<K_{v,q}$,
where the field bounds are in Table~\ref{tab:v-certificates}.  The displayed
$\Sigma_q$ values
give the explicit coefficients $\mathcal C_q=10+3\Sigma_q$; no unprinted density or
endpoint constant is used below.

\subsection{The all-$m$ smoothing estimate}
\label{subsec:v-allm-proof}

\begin{lemma}[logarithmic-potential modulus]
\label{lem:v-potential-modulus}
If a compactly supported positive measure $\eta$ of mass $M$ satisfies
$\eta(B(w,t))\le At$ for every disk, then for $d=|w-w'|\le1/3$,
\[
 |U_\eta(w)-U_\eta(w')|
 \le(6A+2M)d\bigl(1+\log(1/d)\bigr).
\]
\end{lemma}
\begin{proof}
Let $N=\{z:|w-z|\le2d\}$.  On $N$ both distances are at most $3d$.
Approximating the closed disk from outside gives $\eta(N)\le2Ad$.  For
either centre $w_*=w,w'$, layer cake bounds
the integral on $N$ by
\[
 \int_N\log\frac1{|w_*-z|}\,d\eta(z)
 \le\int_0^\infty\min(2Ad,Ae^{-s})\,ds
 =2Ad\bigl(1+\log(1/(2d))\bigr).
\]
Off $N$, the logarithmic difference is at most $2d/|w-z|$, while another
layer-cake integration gives
$\int_{|w-z|>2d}|w-z|^{-1}d\eta(z)\le M+A\log(1/(2d))$.
Adding these bounds gives at most
$6Ad\log(1/d)+2Md$, which is bounded by the displayed expression.
\end{proof}

For $q=2,\ldots,7$, let $\alpha_q$ be the strict lower endpoints in
Table~\ref{tab:v-certificates}.  In particular, the printed table verifies
$\alpha_q>\tau_q$ in each of the six cases.  Then for every
$m=4n\ge64$,
\begin{equation}\label{eq:v-allm-bound}
 \log|G_m(\operatorname{Li}_2(-1/q);-q)|
 \le-\alpha_qm^2+\mathcal C_qm\log m,
 \qquad \mathcal C_q=10+3\Sigma_q.
\end{equation}
\begin{proof}
Fix $q$ and a nonzero integrand point in \eqref{eq:v-integral}.  Put
$\epsilon=m^{-2}$.  Replace each atom $u$ of its empirical measures by
the uniform probability measure on
$u+2i\epsilon+\epsilon e^{i\vartheta}$.  Every point of this circle is
within $3\epsilon$ of $u$.  The common translation preserves pairwise
centres; subharmonicity of $\log|z|$ shows that every off-diagonal positive
logarithmic pair energy does not decrease.  Each diagonal circle has
self-energy $\log\epsilon$, and there are $m+5m/4+m=13m/4$ atoms.  Thus,
writing $\mu^\epsilon$ for the smoothed empirical triple and
$\mathcal Q_{q,\rm disc}$ for the three self-energies with their diagonals
removed,
\[
 \mathcal Q_{q,\rm disc}(\mu)
 \le\mathcal Q_q(\mu^\epsilon)
       +\frac{13}{4m}\log(1/\epsilon)+C_{D,q}\epsilon,
 \quad C_{D,q}=8L_{D,q},
 \quad L_{D,q}=\frac{1+3/4096}{q-(1+3/4096)^2}.
\]
Indeed, on the two smearing disks $|u|,|v|\le1+3\epsilon$ and
$|uv-xv_0|\le6\epsilon+9\epsilon^2$.  The denominator stays at distance
at least $q-(1+3/4096)^2$ from zero.  The mean-value bound for
$\log|q+uv|$, multiplied by the cross mass $5/4$, is at most
\[
\left(\frac{15}{2}+\frac{45}{4}\epsilon\right)L_{D,q}\epsilon
\le C_{D,q}\epsilon,
\]
since $\epsilon\le1/4096$.

The smoothed measures are tame and their $v$-support avoids zero, so
Lemma~\ref{lem:v-concavity} applies.  The modulus lemma and the constants
in Table~\ref{tab:v-transfer-constants} bound the changes in all reference
potentials by
\[
 \Sigma_q(3\epsilon)\bigl(1+\log(1/(3\epsilon))\bigr).
\]
For the two denominator potentials in the dual fields, each derivative in
the field variable is at most $L_{D,q}$ on the tube.  Their total mass
coefficient is $5/2$, so their transfer error is at most
$C_{D,q}\epsilon$.

For the positive one-body terms, if a real centre $t\ge\delta/2$, the
circle mean of $\log|t+2i\epsilon+\epsilon e^{i\vartheta}|$ is
$\log|t+2i\epsilon|\le\log t+8\epsilon^2/\delta^2$.  The combined
weighted loss for $x$ and $v$ is therefore at most
$(17/2)\epsilon^2/\delta^2$.  If $t<\delta/2$, the whole smearing circle
lies in $|w|\le\delta$, so \eqref{eq:v-endpoint} gives the endpoint bound directly.
For the negative $z$ one-body term, $|z|\ge r$,
$|z^\epsilon-z|\le3\epsilon$, and $3\epsilon<r$ give
\[
 \left|\log|z^\epsilon|-\log|z|\right|
 \le\frac{3\epsilon}{r-3\epsilon}.
\]
Consequently,
\[
 -\frac54\log|z^\epsilon|
 \le-\frac54\log|z|+\frac{15\epsilon}{4(r-3\epsilon)}.
\]
The centres lie on the certified supports; combining the last three
estimates with the endpoint bounds and the directed field bounds yields
\begin{align*}
 \frac{\log|\Phi_q|}{m^2}
 \le{}&-\mathcal Q_q(\nu_q)+K_{x,q}+\frac54K_{v,q}+K_{z,q}
 +\frac{13}{4m}\log(1/\epsilon)\\
 &+3\Sigma_q\epsilon\bigl(1+\log(1/(3\epsilon))\bigr)
 +\left(2C_{D,q}+\frac{15}{4(r-3/4096)}\right)\epsilon
 +\frac{17}{2}\frac{\epsilon^2}{\delta^2}.
\end{align*}
Here the coefficient $2C_{D,q}$ accounts separately for the denominator
in the smoothed energy and in the transferred dual fields.  For all
$q\ge2$, the rational comparison
$100(1+3/4096)<101(2-(1+3/4096)^2)$ gives
$L_{D,q}<101/100$, hence $C_{D,q}<8.08$.  Also
$15/(4(r-3/4096))<7$, and
\[
 \frac{17}{2\delta^2m^2}\le\frac{531250}{4096}<130
 \qquad(m\ge64).
\]
The total constant after multiplying by $m^2$ is therefore less than
$16.2+7+130<200$.

The path length is exactly $\mathcal T_q$.  Since $q\le7$ and $e>8/3$, we have
$\log q<2$.  Also $\pi<22/7$, so
$\sqrt{\pi^2+\log^2q}<4$, whence
$\mathcal T_q/q<21/10<e$.  The prefactor in \eqref{eq:v-integral} contributes
$(3m^2/2-m)\log q$, and the path contributes $m\log\mathcal T_q$; their linear
part is $m\log(\mathcal T_q/q)\le m$.  Add these terms and drop the favorable
$-\log((m!)^2L!)$.  With $\epsilon=m^{-2}$, the diagonal term is
$(13/2)m\log m$.  Also
\[
 3\Sigma_q\bigl(1+\log(m^2/3)\bigr)\le3\Sigma_qm\log m,
 \quad 200\le m\log m,
 \quad m\le\tfrac14m\log m
\]
for $m\ge64$.  The resulting coefficient is at most
$13/2+3\Sigma_q+1+1/4<10+3\Sigma_q$.  Finally the definition of $\alpha_q$ in
\eqref{eq:v-alpha} gives \eqref{eq:v-allm-bound}.  If the determinant
integrand vanishes, the bound holds trivially; otherwise all empirical
measures above are defined and the argument applies.  This proves the
claim for every $m=4n\ge64$, not just for the finite certificate values.
\end{proof}

Thus the directed margins and the all-$m$ analytic estimate are separate
inputs: the former supplies $\alpha_q>\tau_q$, while
\eqref{eq:v-allm-bound} supplies the full quantifier with an explicit
finite coefficient for each fixed $q=2,\ldots,7$.
